\section*{Resumen}

Paralegales es una plataforma de \textit{legaltech} que automatiza el seguimiento de procesos judiciales de la Rama Judicial Unificada en Colombia. Su desarrollo inicial, apoyado en el SDK de Firebase para acelerar la salida a producción, dejó deuda técnica relevante: lógica de negocio crítica en el \textit{frontend}, fragmentación arquitectónica entre servicios, ausencia de observabilidad transversal y falta de prácticas formales de aseguramiento de calidad. Descartada la migración total hacia AWS por su costo e impacto sobre un producto ya en operación, este proyecto adopta una estrategia de evolución arquitectónica híbrida en cuatro frentes: reubicación de la lógica crítica hacia el lado del servidor mediante funciones \textit{serverless} y un \textit{Backend-for-Frontend}; refactorización del \textit{frontend} con \textit{Domain-Driven Design} y del \textit{backend} bajo una separación \textit{Handler--Service--Repository}; construcción de un sistema de observabilidad estructurada; y consolidación de prácticas continuas de calidad. El resultado es una plataforma con menor exposición de lógica crítica, mejor trazabilidad operativa y mayor mantenibilidad, preparada para una evolución incremental sin comprometer la continuidad del servicio.

\vspace{1em}

\textbf{Palabras clave:} Arquitectura de software híbrida, \textit{Backend-for-Frontend}, Aseguramiento de calidad, Observabilidad, \textit{Legaltech}.

\newpage

% El documento usa silabeo español; se desactiva la partición de palabras en el
% abstract para evitar cortes incorrectos en el texto en inglés.
\begingroup
\hyphenpenalty=10000
\exhyphenpenalty=10000

\section*{Abstract}

Paralegales is a \textit{legaltech} platform that automates the tracking of judicial proceedings in Colombia's Unified Judicial Branch. Its initial development relied on the Firebase SDK to shorten time to market, which left significant technical debt: critical business logic running in the \textit{frontend}, architectural fragmentation across services, no cross-cutting observability, and no formal quality assurance practices. Since a full migration to AWS was ruled out because of its cost and its impact on a product already in production, this project adopts a hybrid architectural evolution strategy along four fronts: relocating critical business logic to the server side through \textit{serverless} functions and a \textit{Backend-for-Frontend}; refactoring the \textit{frontend} with \textit{Domain-Driven Design} and the \textit{backend} under a \textit{Handler--Service--Repository} separation; building a structured observability system; and consolidating continuous quality assurance practices. The outcome is a platform with less exposure of critical logic, better operational traceability, and improved maintainability, ready for controlled incremental evolution without compromising service continuity.

\vspace{1em}

\textbf{Keywords:} Hybrid software architecture, \textit{Backend-for-Frontend}, Software quality assurance, Observability, \textit{Legaltech}.

\endgroup
